% ---
% id: meat-core/cor:drop-program
% kind: cor
% title: "Convexity, gradient and limit of the dropped program"
% summary: >-
%   The dropped objective is jointly convex and lower semicontinuous, its
%   infimum attained when finite, and on a fully dropped joint sector it reads
%   the block only through its statistics. As the Rényi order tends to one the
%   dropped value tends to the relative-entropy program over the kept pairs.
% prominence: standard
% status: implemented
% uses: [def:dropped, prop:drop-bound, lem:joint-sectors, def:joint-sectors, def:feasible, foundations/lem:kernel, lem:b-derivative, thm:oracle, def:maps, foundations/def:kl, thm:branch-formula, def:admissible-declaration, def:admissible-acceptance-set]
% mentions: [thm:reduction, lem:choi-factor]
% assumes: []
% refs: []
% provenance: newly-derived
% implements: certiqs_key.meat_core.objective.joint_gradient:joint_gradient
% lean:
% notes: "The general gradient formula of the dropped objective on kept blocks is that of cor:full-gradient with the scalar K replaced by the dropped K, by the chain rule on a sum with fewer summands; it is stated in the frank-wolfe submodule where it is consumed."
% ---

The dropped program keeps the undropped program's feasible set and only
removes deficits from each $K$, so the properties the program's numerics
rest on carry over, and a fully dropped joint sector costs no matrix
function. Part~(iii) is written in state coordinates; in ensemble form the
statistics operators are read in the Choi coordinates of
Lemma~\ref{lem:choi-factor}, Alice's blocks not being declared there.

\begin{corollary}\label{cor:drop-program}
Under the hypotheses of Proposition~\ref{prop:drop-bound}, with $\acSet$
convex and compact (Definition~\ref{def:admissible-acceptance-set}):
\begin{enumerate}[label=(\roman*),leftmargin=2.2em]
\item $K^{\mathcal{D}}_{\hat c}$ is concave and continuous on the cone of
operators block diagonal along the joint sectors;
\item $O^{\mathcal{D}}_{\alpha}$ is jointly convex and lower semicontinuous on
$\Delta(\Chatset)\times\Feas$ with values in $[0,+\infty]$, and the infimum
\eqref{eq:hlowD} is attained whenever it is finite;
\item if $(t,s)$ is fully dropped, then $O^{\mathcal{D}}_{\alpha}$ depends on
the block $\omega^{(t,s)}$ only through the real numbers
$\tr[(M^{(t)}_x\otimes F^{(s)}_v)\,\omega^{(t,s)}]$, $(x,v)\in\Xset\times\Vset$;
at a point where every $P_{\hat c}$ with $q(\hat c)>0$ and every
$K^{\mathcal{D}}_{\hat c}$ with $q(\hat c)>0$, $\hat c\in\Chkey$, is
positive, its derivative in the direction of a Hermitian $\Delta$ supported
on that joint sector is
\begin{equation}\label{eq:drop-gradient}
  -\log_2\!\mathrm{e}\sum_{\hat c\in\Chatset}\beta_{\hat c}\,
  \tr\big[\Gamma^{(t,s)}_{\hat c}\Delta\big],\qquad
  \Gamma^{(t,s)}_{\hat c}\coloneqq\!\!\sum_{(x,v):\,\phi(a(x,v))=\hat c}\!\!
  M^{(t)}_x\otimes F^{(s)}_v,\qquad
  \beta_{\hat c}\coloneqq
  \begin{cases}
   \frac{\alpha}{\alpha-1}\frac{q(\hat c)}{P_{\hat c}},&\hat c\notin\Chkey,\\[0.6ex]
   \frac{q(\hat c)}{P_{\hat c}}+\frac{1}{\alpha-1}
   \frac{q(\hat c)}{K^{\mathcal{D}}_{\hat c}},&\hat c\in\Chkey,
  \end{cases}
\end{equation}
with $\beta_{\hat c}\coloneqq0$ where $q(\hat c)=0$;
\item under the additional hypotheses of Theorem~\ref{thm:oracle},
\begin{equation}\label{eq:drop-oracle}
  \lim_{\alpha\downarrow1}\underline{h}^{\mathcal{D}}_{\alpha}
  =\inf_{\omega\in\Feas_1}\ \sum_{\hat c\in\Chkey}\
  \sum_{\substack{(t,s)\in\mathbb{J}_K(\hat c)\\(\hat c,(t,s))\notin\mathcal{D}}}
  \Dtwo\Big(\cG_{\hat c,(t,s)}\big(\omega^{(t,s)}\big)\,\Big\|\,
  \cZ\big(\cG_{\hat c,(t,s)}\big(\omega^{(t,s)}\big)\big)\Big).
\end{equation}
\end{enumerate}
\end{corollary}

\begin{proof}
\emph{Idea.} Each $K^{\mathcal{D}}_{\hat c}$ is $P_{\hat c}$ plus a sum of
concave deficits, so (i) holds summand by summand, and (ii) follows from the
monotonicity and joint convexity of the kernel as for any concave,
continuous, non-negative argument, with attainment from compactness. On a
fully dropped joint sector only the
linear $P_{\hat c}$ read the block, which gives (iii). For (iv) the proof of
Theorem~\ref{thm:oracle} goes through with the kept summands in place of
all of them; we list the changes.

(i) $P_{\hat c}$ is linear. Each $f_{\hat c,(t,s)}$ is concave on its cone
(Lemma~\ref{lem:joint-sectors}\,(ii)), and
$\omega\mapsto\omega^{(t,s)}$ is linear, so \eqref{eq:KD} is a sum of concave
functions. Continuity: $X\mapsto X^{1/\alpha}$ and $Y\mapsto Y^{\alpha}$ are
continuous on the positive semidefinite cone by the continuous functional
calculus, $\cZ$, $\cG_{\hat c,s}$ and the trace are linear, so each
$f_{\hat c,(t,s)}$ of \eqref{eq:f-def} is continuous.

(ii) Write $O^{\mathcal{D}}_{\alpha}$ as the sum over $\hat c$ of
$\frac{\alpha}{\alpha-1}d(q(\hat c),P_{\hat c}(\omega))$ for
$\hat c\notin\Chkey$ and of
$d(q(\hat c),P_{\hat c}(\omega))+\frac{1}{\alpha-1}
d(q(\hat c),K^{\mathcal{D}}_{\hat c}(\omega))$ for $\hat c\in\Chkey$, with
$d$ the kernel of \xref{foundations/lem:kernel}. Let $g$ be concave, continuous and
non-negative on a convex set, as $P_{\hat c}$ and, by (i) and
Proposition~\ref{prop:drop-bound}\,(i), $K^{\mathcal{D}}_{\hat c}$ are on the
convex set $\Feas$. For a convex combination with weights $\lambda$,
$1-\lambda$: $g(\omega_{\lambda})\ge\lambda g(\omega_1)+(1-\lambda)g(\omega_2)$,
so, $d$ being non-increasing in its second argument and jointly convex,
$d(q_{\lambda},g(\omega_{\lambda}))\le
d\big(q_{\lambda},\lambda g(\omega_1)+(1-\lambda)g(\omega_2)\big)\le
\lambda d(q_1,g(\omega_1))+(1-\lambda)d(q_2,g(\omega_2))$. For lower
semicontinuity let $(q_n,\omega_n)\to(q,\omega)$; then
$g(\omega_n)\to g(\omega)$ and lower semicontinuity of $d$ gives
$\liminf_nd(q_n,g(\omega_n))\ge d(q,g(\omega))$. Sums with positive
coefficients preserve both properties. Non-negativity is
Proposition~\ref{prop:drop-bound}\,(ii). $\acSet\times\Feas$ is compact
(Definition~\ref{def:admissible-acceptance-set}, Definition~\ref{def:feasible}), and a lower
semicontinuous function with values in $(-\infty,+\infty]$ attains its
infimum on a non-empty compact set.

(iii) If $(t,s)$ is fully dropped, no $f_{\hat c,(t,s)}$ occurs in any
$K^{\mathcal{D}}_{\hat c}$, so the block enters
$O^{\mathcal{D}}_{\alpha}$ through the values $P_{\hat c}(\omega)$ alone,
and its contribution to $P_{\hat c}$ is
$P_{\hat c,(t,s)}(\omega^{(t,s)})
=\tr[\Gamma^{(t,s)}_{\hat c}\omega^{(t,s)}]$ by \eqref{eq:Phimap}, a sum of
the stated numbers. With every relevant argument positive,
$\frac{\partial}{\partial v}\,u\log_2\frac uv=-\log_2\!\mathrm{e}\,\frac uv$,
and the directional derivative of $P_{\hat c}$ and, the deficits of the
block being absent, of $K^{\mathcal{D}}_{\hat c}$ along $\Delta$ is
$\tr[\Gamma^{(t,s)}_{\hat c}\Delta]$ in both cases. Collecting the prefactors
of (ii) gives \eqref{eq:drop-gradient}; terms with $q(\hat c)=0$ are
identically zero.

(iv) We list the changes to the proof of Theorem~\ref{thm:oracle}, which is
written for $\mathcal{D}=\varnothing$ and the index $s$. Write
$\sigma_{\hat c,(t,s)}(\omega)\coloneqq\cG_{\hat c,(t,s)}(\omega^{(t,s)})$,
an operator of trace at most $1$ for $\omega\in\Feas$, and
$e^{\mathcal{D}}_{\hat c}\coloneqq\frac{1}{1-\alpha}\log_2
\frac{K^{\mathcal{D}}_{\hat c}}{P_{\hat c}}$, which is $\ge0$ by
Proposition~\ref{prop:drop-bound}\,(i). By \eqref{eq:KD}, \eqref{eq:f-def}
and $B_{\sigma}(1)=\tr\sigma$,
\begin{equation*}
  K^{\mathcal{D}}_{\hat c}(\omega)-P_{\hat c}(\omega)
  =\sum_{(t,s):\,(\hat c,(t,s))\notin\mathcal{D}}
  \Big(B_{\sigma_{\hat c,(t,s)}(\omega)}(\alpha)
  -B_{\sigma_{\hat c,(t,s)}(\omega)}(1)\Big),
\end{equation*}
which replaces the identity for $K_{\hat c}-P_{\hat c}$ used there.
\emph{Upper bound:} with $\mathbf q=\Phimap[\omega]$ the divergence terms
vanish and $O^{\mathcal{D}}_{\alpha}(\mathbf q,\omega)
=\sum_{\hat c\in\Chkey}P_{\hat c}\,e^{\mathcal{D}}_{\hat c}$; the limit is
taken summand by summand with Lemma~\ref{lem:b-derivative}, exactly as there,
and gives the sum of relative entropies over the kept pairs; pairs that are
key-inert contribute zero because their $\sigma$ is block diagonal for
$\cZ$ (Lemma~\ref{lem:joint-sectors}\,(iii) and its proof), which restricts
the sum to $\mathbb{J}_K(\hat c)$. \emph{Lower bound:} the decomposition of
the objective into the divergence part and
$\sum_{\hat c}q(\hat c)e^{\mathcal{D}}_{\hat c}$ is the identity
\eqref{eq:rearrange} of real numbers, valid for any $K\in(0,P]$; both parts
are non-negative; compactness of $\acSet\times\Feas$, lower semicontinuity of
the divergence and the conclusion $\omega^{*}\in\Feas_1$ do not involve $K$;
$u_n\coloneqq(K^{\mathcal{D}}_{\hat c}-P_{\hat c})/P_{\hat c}\in[-1,0]$ by
Proposition~\ref{prop:drop-bound}\,(i); and the mean-value step with the
majorant \eqref{eq:majorant} and the joint continuity of $c_{\sigma}(\alpha)$
is applied to each kept summand, the maps
$\omega\mapsto\sigma_{\hat c,(t,s)}(\omega)$ being continuous. The resulting
$\liminf$ is the right-hand side of \eqref{eq:drop-oracle} at $\omega^{*}$.
\end{proof}
